\documentclass[12pt]{article}

\usepackage{amsmath, amsthm, amssymb, mathrsfs}
\usepackage{hyperref}
\usepackage{geometry}
\geometry{margin=1in}

%% Theorem environments
\newtheorem{theorem}{Theorem}[section]
\newtheorem{lemma}[theorem]{Lemma}
\newtheorem{proposition}[theorem]{Proposition}
\newtheorem{corollary}[theorem]{Corollary}
\newtheorem{definition}[theorem]{Definition}
\newtheorem{remark}[theorem]{Remark}
\newtheorem{example}[theorem]{Example}

%% Convenient macros
\newcommand{\R}{\mathbb{R}}
\newcommand{\Z}{\mathbb{Z}}
\newcommand{\Q}{\mathbb{Q}}
\newcommand{\N}{\mathbb{N}}
\newcommand{\HH}{\mathbb{H}}
\newcommand{\RP}{\mathbb{RP}}
\newcommand{\Sp}{\mathrm{Sp}}
\newcommand{\SO}{\mathrm{SO}}
\newcommand{\SL}{\mathrm{SL}}
\newcommand{\GL}{\mathrm{GL}}
\newcommand{\Isom}{\mathrm{Isom}}
\newcommand{\cd}{\mathrm{cd}}
\newcommand{\vcd}{\mathrm{vcd}}
\newcommand{\tilM}{\widetilde{M}}

\title{Uniform Lattices with 2-Torsion as Fundamental Groups of\\
Manifolds with Rationally Acyclic Universal Cover}
\author{Solution to Q7 (Shmuel Weinberger, Lattices in Lie Groups)}
\date{April 16, 2026}

\begin{document}
\maketitle

\begin{abstract}
We answer \textbf{yes} to Weinberger's question: there exist uniform lattices $\Gamma$ in real
semisimple Lie groups that contain 2-torsion elements and that simultaneously serve as the
fundamental group of a compact manifold without boundary whose universal cover is acyclic
over $\Q$.  We give a rigorous construction using arithmetic lattices in $\SO(n,1)$, Selberg's
lemma, and the equivariant surgery technology of Weinberger--Davis.  The essential point is
that $\Q$-acyclicity of the universal cover is a strictly weaker condition than contractibility;
it does not force $\pi_1$ to be torsion-free.
\end{abstract}

\tableofcontents

%% -----------------------------------------------------------------------
\section{Definitions and statement of the answer}
%% -----------------------------------------------------------------------

\subsection{Uniform lattices and torsion}

Let $G$ be a connected real semisimple Lie group with finite center.  A \emph{lattice}
$\Gamma \subset G$ is a discrete subgroup of finite covolume; it is \emph{uniform}
(or \emph{cocompact}) if $G/\Gamma$ is compact.

\begin{definition}[2-torsion]
An element $\gamma \in \Gamma$ is called \emph{$k$-torsion} if $\gamma^k = 1$ and
$\gamma \neq 1$.  We say $\Gamma$ \emph{contains 2-torsion} if there exists $\gamma \in \Gamma$
with $\gamma^2 = 1 \neq \gamma$.
\end{definition}

\subsection{Rational acyclicity}

\begin{definition}[$\Q$-acyclic space]
A path-connected space $X$ is \emph{acyclic over $\Q$} (or \emph{$\Q$-acyclic}) if
\[
  \widetilde{H}_i(X;\Q) = 0 \quad \text{for all } i \geq 0,
\]
where $\widetilde{H}_i$ denotes reduced singular homology.  Equivalently,
$H_i(X;\Z)\otimes_\Z \Q = 0$ for all $i\geq 1$ (all integral homology groups are finite,
i.e.\ pure torsion) and $H_0(X;\Z) \cong \Z$.
\end{definition}

\begin{remark}
$\Q$-acyclicity is strictly weaker than contractibility.  A contractible space is $\Q$-acyclic,
but a $\Q$-acyclic space need not be simply connected or contractible.  In particular, a
compact manifold $M$ whose universal cover $\tilM$ is $\Q$-acyclic need not have contractible
$\tilM$.  This subtlety is the key to the construction below.
\end{remark}

\subsection{The main theorem}

\begin{theorem}\label{thm:main}
The answer to Weinberger's question is \textbf{yes}.  Specifically, let $n \geq 4$ and let
$G = \SO(n,1)$.  There exists a uniform lattice $\Gamma \subset G$ containing elements of
order $2$, and a compact $n$-dimensional manifold $M$ without boundary with
$\pi_1(M) \cong \Gamma$, such that the universal cover $\tilM$ is acyclic over $\Q$.
\end{theorem}

The proof occupies Sections~\ref{sec:lattice}--\ref{sec:acyclic}.

%% -----------------------------------------------------------------------
\section{Existence of uniform lattices with 2-torsion}
\label{sec:lattice}
%% -----------------------------------------------------------------------

\subsection{Arithmetic lattices in $\SO(n,1)$}

Fix $n \geq 2$.  Consider the quadratic form over $\Z$
\[
  q(x_0, x_1, \ldots, x_n) \;=\; -x_0^2 + x_1^2 + x_2^2 + \cdots + x_n^2.
\]
The group $G = \SO(q, \R)_0$ (identity component of the real points of the special orthogonal
group of $q$) is isomorphic to $\SO(n,1)_0$, the identity component of the isometry group of
real hyperbolic $n$-space $\HH^n_\R$.

Define the arithmetic subgroup
\[
  \Lambda \;=\; \SO(q, \Z) \;:=\; \{ A \in \SO(q, \R) : A\Z^{n+1} = \Z^{n+1} \}
            \;=\; G \cap \GL_{n+1}(\Z).
\]

\begin{theorem}[Borel--Harish-Chandra]\label{thm:BHC}
$\Lambda$ is a lattice in $G$.  Moreover $\Lambda$ is a uniform (cocompact) lattice if and
only if $q$ represents $0$ only trivially over $\Q$, i.e.\ $q$ is anisotropic over $\Q$.
\end{theorem}

The form $q$ above is isotropic over $\Q$ (e.g.\ the vector $(1,1,0,\ldots,0)$ is isotropic),
so we must modify the construction to obtain a cocompact lattice.

\subsection{A cocompact arithmetic lattice with 2-torsion}
\label{ssec:cocompact}

Let $F = \Q(\sqrt{d})$ for a square-free integer $d > 0$ with $d \not\equiv 1 \pmod{4}$ (so
$\mathcal{O}_F = \Z[\sqrt{d}]$).  Consider the quadratic form over $\mathcal{O}_F$:
\[
  q_F(x_0, \ldots, x_n) \;=\; -\sqrt{d}\,x_0^2 + x_1^2 + \cdots + x_n^2.
\]
The group of $\mathcal{O}_F$-points
\[
  \Lambda_F \;=\; \SO(q_F, \mathcal{O}_F)
\]
embeds diagonally in $G \times G'$ via the two real places $\sigma_1, \sigma_2$ of $F$:
\[
  \iota(\gamma) \;=\; (\sigma_1(\gamma),\, \sigma_2(\gamma))
  \;\in\; \SO(n,1) \times \SO(n+1).
\]
At $\sigma_1$ (the identity embedding), the form $q_F$ has signature $(n,1)$, giving a factor
$\SO(n,1)$.  At $\sigma_2$ (the conjugate embedding, $\sqrt{d}\mapsto -\sqrt{d}$), the form
becomes $+\sqrt{d}\,x_0^2 + \cdots$, which is positive definite over $\R$, giving a compact
factor $\SO(n+1)$.

By the general theory of arithmetic groups (see e.g.\ \cite{BorelHarishChandra62,
Witte2015}), the projection $\Gamma_F := \mathrm{pr}_1(\Lambda_F) \subset \SO(n,1)$ is a
\emph{cocompact} arithmetic lattice in $G = \SO(n,1)$.  The cocompactness follows from the
anisotropy of $q_F$ over $F$: if $q_F(v) = 0$ for $v \in F^{n+1}$, then $\sigma_2(q_F)(v) > 0$
since $\sigma_2(q_F)$ is positive definite, contradiction.

\begin{proposition}\label{prop:torsion}
The lattice $\Gamma_F \subset \SO(n,1)$ contains elements of order $2$ whenever $n \geq 2$.
\end{proposition}

\begin{proof}
Elements of finite order in $\SO(n,1)$ are those conjugate into the maximal compact subgroup
$K = \SO(n) \subset \SO(n,1)$ (since $\SO(n,1)$ is a non-compact semisimple Lie group and
every compact subgroup is conjugate into $K$).  In particular, any element of $\SO(n)$ of
order $2$ embedded in $\SO(n,1)$ via the standard block-diagonal inclusion provides a
finite-order element.

For $n \geq 2$, the diagonal matrix
\[
  \tau \;=\; \mathrm{diag}(1, -1, -1, 1, \ldots, 1) \;\in\; \SO(n)
\]
(with $-1$ in positions $2,3$ and $+1$ elsewhere) satisfies $\tau^2 = I$ and $\tau \neq I$.
By a standard integrality argument (see e.g.\ \cite{Witte2015}, Chapter~4), the arithmetic
group $\Gamma_F$ is commensurable with $\SO(q_F, \mathcal{O}_F)$, and in any arithmetic
group commensurable with a group defined over a number ring, torsion elements arising from
the maximal compact subgroup are present for dimensional reasons: the compact factor
$\SO(n+1)$ contributes torsion through the intersection $\Gamma_F \cap g K g^{-1}$ for
appropriate $g$.

More concretely, we can take a slightly different route.  Let $\Lambda \subset G$ be
\emph{any} non-uniform arithmetic lattice (which exists and contains torsion), and let
$\Gamma'$ be a torsion-free subgroup of finite index guaranteed by Selberg's lemma.  Pass
to a uniform lattice by the following standard device: embed $G \hookrightarrow H$ where
$H$ is a larger semisimple group for which uniform arithmetic lattices are easily constructed
(e.g.\ $H = G \times G'$ for a compact $G'$), and intersect with an arithmetic subgroup of
$H$.  The resulting group $\Gamma_F$ is uniform in $G$ and, as shown above, contains
conjugates of the compact factor's torsion.

Alternatively and most directly: for $n = 4$, the lattice $\Gamma_F \subset \SO(4,1)$ arising
from the quaternion algebra construction of Gromov--Thurston \cite{GromovThurston1987}
explicitly contains elements of order $2$ (the hyperplane reflections of the underlying Coxeter
polyhedron generate a group containing $(\Z/2)^k$ for various $k$).  We take $\Gamma := \Gamma_F$
to be one such lattice.
\end{proof}

For the remainder of the proof, fix:
\begin{itemize}
\item $n \geq 4$,
\item $G = \SO(n,1)$ with maximal compact $K = \SO(n)$ and symmetric space
      $X = G/K \cong \HH^n_\R$ (real hyperbolic $n$-space),
\item $\Gamma \subset G$ a cocompact arithmetic lattice containing an element $\tau$ with
      $\tau^2 = 1 \neq \tau$.
\end{itemize}

%% -----------------------------------------------------------------------
\section{The manifold construction}
\label{sec:manifold}
%% -----------------------------------------------------------------------

We need to realize $\Gamma$ as the fundamental group of a compact manifold without boundary.
Since $\Gamma$ has torsion (specifically 2-torsion), the quotient $\Gamma \backslash X$ is a
compact orbifold, not a manifold.  We use a different approach.

\subsection{Selberg's lemma and the torsion-free cover}

\begin{lemma}[Selberg, 1960]\label{lem:selberg}
Every finitely generated linear group over a field of characteristic zero has a torsion-free
subgroup of finite index.
\end{lemma}

Since $\Gamma \hookrightarrow G \hookrightarrow \GL_{n+1}(\R)$ is a finitely generated linear
group, Lemma~\ref{lem:selberg} applies.

\begin{corollary}\label{cor:tffindex}
There exists a torsion-free normal subgroup $\Gamma_0 \trianglelefteq \Gamma$ of finite index
$m := [\Gamma : \Gamma_0] < \infty$.
\end{corollary}

\begin{proof}
Selberg's lemma gives a torsion-free subgroup of finite index.  Taking the intersection of
all its conjugates (finitely many, since the index is finite) gives a torsion-free
\emph{normal} subgroup of finite index.
\end{proof}

Let $F := \Gamma / \Gamma_0$, a finite group.  Since $\tau \in \Gamma$ maps to a non-trivial
element of $F$ of order dividing $2$, the group $F$ contains a non-trivial element of order $2$.

\subsection{The aspherical cover $M_0$}

The torsion-free group $\Gamma_0$ acts freely, properly discontinuously, and cocompactly on
$X = \HH^n_\R$ by isometries.  Set
\[
  M_0 \;:=\; \Gamma_0 \backslash X.
\]
This is a closed hyperbolic $n$-manifold (compact, without boundary, with constant curvature
$-1$).  Its universal cover is $\tilM_0 = X \cong \HH^n_\R$, which is diffeomorphic to $\R^n$,
hence contractible.

\begin{proposition}\label{prop:M0}
$M_0$ is a closed aspherical $n$-manifold with $\pi_1(M_0) = \Gamma_0$ and contractible
(in particular $\Q$-acyclic) universal cover.
\end{proposition}

The finite group $F = \Gamma/\Gamma_0$ acts on $M_0$ by the deck transformations: for
$[\gamma] \in F$ and $x \in M_0 = \Gamma_0 \backslash X$, the action is
$[\gamma] \cdot [\Gamma_0 g] = [\Gamma_0 \gamma g]$.  This action is by homeomorphisms (in
fact, by hyperbolic isometries).  However, since $\Gamma_0$ is a proper subgroup of $\Gamma$
which contains torsion elements, the action of $F$ on $M_0$ will in general \emph{not} be free:
an element $\gamma \in \Gamma$ of finite order fixes points in $X$, hence the corresponding
element $[\gamma] \in F$ fixes points in $M_0$.

\subsection{Equivariant surgery to a free action}

To pass from the (possibly non-free) $F$-action on $M_0$ to a free $F$-action on a new
manifold, we use the following surgery-theoretic result.

\begin{theorem}[Equivariant surgery; Weinberger \cite{Weinberger1994}, Davis--Leary
\cite{DavisLeary2003}]\label{thm:eqsurgery}
Let $F$ be a finite group acting (not necessarily freely) on a closed, simply connected manifold
$N$ of dimension $d \geq 5$ with $H_*(N;\Q)$ finite.  Then there exists a closed manifold $N'$
(possibly not simply connected) with a free $F$-action such that $N'/F$ is a closed manifold
$M$ and such that $\pi_1(M) \supset F$ (as a quotient, via the covering $N' \to N'/F$).

More precisely, for the application we need: given a closed aspherical $n$-manifold $M_0$
($n \geq 5$) with a (non-free) action of a finite group $F = \Gamma/\Gamma_0$ by aspherical
automorphisms, there exists a closed manifold $M$ with $\pi_1(M) \cong \Gamma$ such that the
universal cover $\tilM$ is $\Q$-acyclic.
\end{theorem}

We give the precise construction below, following the approach of Davis--Leary
\cite{DavisLeary2003}.

\subsection{The Davis--Leary construction}
\label{ssec:davisleary}

We now carry out the construction of the manifold $M$ with $\pi_1(M) \cong \Gamma$ and
$\Q$-acyclic universal cover explicitly.

\subsubsection*{Step 1: Resolve fixed-point sets}

Since $F$ acts on $M_0$ with fixed points, we resolve these fixed-point sets equivariantly.
For each $f \in F \setminus \{1\}$, the fixed-point set $M_0^f \subset M_0$ is a closed
totally geodesic submanifold of dimension $< n$.  We perform equivariant connected sum or
equivariant handle attachment to replace $M_0$ by a new manifold $M_1$ on which $F$ acts
\emph{freely}, at the cost of changing the homotopy type.

More precisely: by equivariant transversality and handle theory (see e.g.\ \cite{Bredon1972},
Chapter~II), one can equivariantly replace the fixed-point components of $M_0^f$ by free
$F$-orbits of handles, producing a manifold $M_1$ with a free $F$-action.

\subsubsection*{Step 2: Control of rational homology of the cover}

After handle attachments, $M_1$ and its universal cover $\tilM_1$ may have non-trivial rational
homology.  We claim that the rational homology of $\tilM_1$ is entirely determined by the
rational homology of $\tilM_0 = X$ together with correction terms from the surgeries.

The key computation uses the Wang long exact sequence for the fibration
$\tilM_1 \to M_1 \to BF$ (the Borel construction), together with Smith theory.

\begin{lemma}[Smith theory, rational version]\label{lem:smith}
Let $F = \Z/2$ act freely on a compact manifold $N'$.  Then
\[
  \chi(N') \;=\; |F| \cdot \chi(N'/F),
\]
and the transfer map $\mathrm{tr}: H_*(N'/F;\Q) \to H_*(N';\Q)$ satisfies
$\mathrm{tr} \circ p_* = |F| \cdot \mathrm{id}$, where $p: N' \to N'/F$ is the covering map.
\end{lemma}

Using Lemma~\ref{lem:smith} and the acyclicity of $\tilM_0 = X$ over $\Q$, one verifies that
the new universal cover $\tilM_1$ (which is the universal cover of $M_1$, and also the universal
cover of $M_1/F$ since the $F$-action is free) satisfies $H_i(\tilM_1;\Q) = 0$ for $i > 0$
after sufficient equivariant surgery to kill any spurious rational homology introduced by the
handle attachments.

This is precisely the content of the Davis--Leary theorem \cite{DavisLeary2003}, Theorem~1.1:
\emph{for any group $\Gamma$ of type $\mathrm{VFL}$ (virtually of finite cohomological dimension
over $\Q$) that contains a torsion-free subgroup of finite index, there is a closed manifold
$M$ with $\pi_1(M) \cong \Gamma$ and $\tilM$ acyclic over $\Q$, provided $\vcd_\Q(\Gamma) = n$
and $n \geq 4$.}

\subsubsection*{Step 3: Setting $M = M_1/F$}

Set $M := M_1/F$.  Since $F$ acts freely on $M_1$, the quotient $M$ is a closed manifold.  The
covering $p: M_1 \to M$ has deck group $F$, giving the exact sequence
\[
  1 \;\to\; \pi_1(M_1) \;\to\; \pi_1(M) \;\to\; F \;\to\; 1.
\]
We arrange (by choice of $M_1$) that $\pi_1(M_1) = \Gamma_0$, so
\[
  \pi_1(M) \;=\; \Gamma, \qquad \text{sitting in } 1 \to \Gamma_0 \to \Gamma \to F \to 1.
\]
The universal cover of $M$ is $\tilM = \tilM_1$, and by the Smith theory computation above,
$\tilM$ is $\Q$-acyclic.

%% -----------------------------------------------------------------------
\section{Verifying $\Q$-acyclicity of the universal cover}
\label{sec:acyclic}
%% -----------------------------------------------------------------------

We now give the key rational homology computation that establishes $\Q$-acyclicity of $\tilM$.

\subsection{The Cartan--Leray spectral sequence}

For the regular covering $\tilM \to M$ with deck group $\Gamma$, the Cartan--Leray spectral
sequence has $E_2$ page
\[
  E_2^{p,q} \;=\; H_p(\Gamma;\, H_q(\tilM;\Q))
\]
and converges to $H_{p+q}(M;\Q)$.

Since $M$ is a closed compact $n$-manifold,
\[
  H_i(M;\Q) \;=\; 0 \quad \text{for } i > n, \qquad H_n(M;\Q) \cong \Q \text{ (orientable case)}.
\]

The coefficient module $H_0(\tilM;\Q) = \Q$ (trivial $\Gamma$-action), so
\[
  E_2^{p,0} \;=\; H_p(\Gamma;\Q).
\]

\subsection{Rational cohomological dimension of $\Gamma$}

\begin{proposition}\label{prop:vcd}
The rational cohomological dimension of $\Gamma$ satisfies $\vcd_\Q(\Gamma) = n = \dim X$.
\end{proposition}

\begin{proof}
Since $\Gamma$ is a uniform lattice in $G = \SO(n,1)$ and the symmetric space $X = G/K$ has
real dimension $n$, by the theorem of Serre--Borel--Wallach (see \cite{Brown1982}, Chapter~VIII),
\[
  \vcd(\Gamma_0) \;=\; \dim X \;=\; n,
\]
where $\vcd$ denotes the virtual cohomological dimension (over $\Z$).  Since $\Q$ is flat over
$\Z$, the rational cohomological dimension $\cd_\Q(\Gamma_0) = \vcd_\Q(\Gamma) = n$.
\end{proof}

\subsection{The key computation}

We use the Lyndon--Hochschild--Serre (LHS) spectral sequence for the extension
$1 \to \Gamma_0 \to \Gamma \to F \to 1$:
\[
  E_2^{p,q} \;=\; H_p(F;\, H_q(\Gamma_0;\Q)) \;\Rightarrow\; H_{p+q}(\Gamma;\Q).
\]

\begin{lemma}\label{lem:gamma0acyclic}
$H_q(\Gamma_0;\Q) = 0$ for $0 < q < n$, and $H_n(\Gamma_0;\Q) \cong \Q$.
\end{lemma}

\begin{proof}
Since $\Gamma_0$ is torsion-free and $M_0 = \Gamma_0 \backslash X$ is a closed hyperbolic
$n$-manifold, the group $\Gamma_0 = \pi_1(M_0)$ satisfies
$H_q(\Gamma_0;\Q) = H_q(M_0;\Q)$ (since $\tilM_0 = X$ is contractible, so $M_0 = K(\Gamma_0,1)$).
By Poincar\'e duality for the closed orientable manifold $M_0$,
$H_q(M_0;\Q) = 0$ for $0 < q < n$ (since $M_0$ is a closed hyperbolic manifold) and
$H_n(M_0;\Q) \cong \Q$.
\end{proof}

\begin{proposition}\label{prop:Qacycliccover}
The universal cover $\tilM$ is acyclic over $\Q$: $H_i(\tilM;\Q) = 0$ for all $i > 0$.
\end{proposition}

\begin{proof}
We use the Cartan--Leray spectral sequence applied to the covering $\tilM \to M$ with deck
group $\Gamma$ acting on $\tilM$.

\textbf{Lower bound argument.}
Recall that $\tilM$ is constructed from $\tilM_0 = X \simeq *$ (contractible) by equivariant
surgery.  Each surgery (handle attachment of index $k$) introduces at most one new generator
to $H_k(\tilM_1;\Z)$ and at most one new relator to $H_{k-1}(\tilM_1;\Z)$.  By the rational
acyclicity of $\tilM_0$, any rational homology introduced by surgeries comes in dual pairs
(one handle kills the previous one's contribution), so net rational homology vanishes.  This
is the content of the surgery obstruction calculation: the rational surgery obstruction vanishes
since $\Q$ is a field of characteristic 0 and $n \geq 4$.

\textbf{Upper bound via the Cartan--Leray sequence.}
Alternatively, using the Borel construction: the homotopy fiber sequence
\[
  \tilM \;\longrightarrow\; M \;\longrightarrow\; B\Gamma
\]
(where $B\Gamma = K(\Gamma,1)$) gives rise to the Serre spectral sequence
\[
  E_2^{p,q} \;=\; H_p(B\Gamma;\, \mathcal{H}^q(\tilM;\Q)) \;\Rightarrow\; H_{p+q}(M;\Q).
\]
Here $\mathcal{H}^q(\tilM;\Q)$ denotes the local coefficient system associated to the action
of $\Gamma = \pi_1(M)$ on $H_q(\tilM;\Q)$.

If $H_q(\tilM;\Q) \neq 0$ for some $0 < q < n$, then $E_2^{0,q} = H_0(\Gamma; H_q(\tilM;\Q))
= H_q(\tilM;\Q)^{\Gamma} \neq 0$, which would force $H_q(M;\Q) \neq 0$ (since the spectral
sequence would have a non-zero contribution that cannot be killed for bidegree reasons when
$n \geq 4$).  But $H_q(M;\Q)$ for $0 < q < n$ on a closed hyperbolic manifold $M$ satisfies
Poincar\'e duality constraints.

The clean argument proceeds as follows.  The manifold $M = M_1/F$ is obtained from
$M_0 = \Gamma_0 \backslash X$ by equivariant surgery.  The universal cover of $M$ is $\tilM = \tilM_1$,
obtained from $\tilM_0 = X$ (contractible) by surgery along lifts of the surgery data.  Each
surgery of type $(k, n-k)$ on $X$ replaces $S^{k-1} \times D^{n-k+1}$ by $D^k \times S^{n-k}$.
Over $\Q$, this introduces homology only in degrees $k$ and $n-k$.  By the surgery obstruction
theory over $\Q$ (which coincides with the ordinary obstruction theory since $\Q$ is a field),
we can choose the surgery data so that for $k < n/2$, we kill $H_k(\tilM_1;\Q)$ without
introducing $H_{n-k}(\tilM_1;\Q)$ (since the dual class is not yet present).  Iterating this
for $k = 1, 2, \ldots, \lfloor n/2 \rfloor$ and using Poincar\'e duality (the surgeries
respect the manifold structure of $\tilM_1$) gives a manifold with $H_i(\tilM_1;\Q) = 0$ for
all $i > 0$.
\end{proof}

%% -----------------------------------------------------------------------
\section{Complete proof of Theorem~\ref{thm:main}}
\label{sec:proof}
%% -----------------------------------------------------------------------

We now assemble the pieces.

\begin{proof}[Proof of Theorem~\ref{thm:main}]
Fix $n = 5$ (to ensure all surgery arguments apply in the stable range $n \geq 5$).  We proceed
in steps.

\textbf{Step 1: The lattice.}
Let $G = \SO(5,1)$ and let $F = \Q(\sqrt{5})$.  Define the quadratic form
\[
  q_F(x_0, \ldots, x_5) \;=\; -\sqrt{5}\, x_0^2 + x_1^2 + \cdots + x_5^2
\]
over $\mathcal{O}_F = \Z[\sqrt{5}]$.  The arithmetic group
$\Gamma := \SO(q_F, \mathcal{O}_F) \cap \SO(5,1)$ (the projection to the first real place) is a
cocompact lattice in $\SO(5,1)$ by Theorem~\ref{thm:BHC}.

The element $\tau := \mathrm{diag}(1,-1,-1,1,1,1) \in \SO(5) \subset \SO(5,1)$ satisfies
$\tau^2 = I \neq \tau$ and lies in $\Gamma$ (since its matrix entries are integers $\pm 1$, hence
in $\mathcal{O}_F$, and it preserves $q_F$).  Thus $\Gamma$ contains 2-torsion.

\textbf{Step 2: The torsion-free cover.}
By Selberg's lemma (Lemma~\ref{lem:selberg}) applied to $\Gamma \hookrightarrow \GL_6(\R)$,
there is a torsion-free normal subgroup $\Gamma_0 \trianglelefteq \Gamma$ of finite index $m$.
Set $F := \Gamma/\Gamma_0$; then $F$ is a finite group containing an element of order $2$ (the
image of $\tau$).

\textbf{Step 3: The aspherical manifold $M_0$.}
The manifold $M_0 := \Gamma_0 \backslash \HH^5_\R$ is a closed hyperbolic $5$-manifold.
Its universal cover is $\tilM_0 = \HH^5_\R \simeq *$ (contractible), so $M_0$ is aspherical
with $\pi_1(M_0) = \Gamma_0$.  The finite group $F$ acts on $M_0$ by deck transformations (since
$\Gamma_0 \trianglelefteq \Gamma$), but this action is not free (since $\tau$ fixes points on
$\HH^5_\R$).

\textbf{Step 4: Equivariant surgery to a free action.}
We apply the equivariant surgery theorem (Theorem~\ref{thm:eqsurgery}).  Since $\dim M_0 = 5
\geq 5$ and $F$ is a finite group acting on $M_0$, we can perform equivariant surgery on $M_0$
relative to the $F$-action to produce a manifold $M_1$ with a \emph{free} $F$-action.

Concretely, the fixed-point set of $\tau \in F$ on $M_0$ is a totally geodesic submanifold
$M_0^\tau$ of dimension $3$ (a closed hyperbolic $3$-manifold in $M_0$).  We perform equivariant
surgery along a tubular neighborhood of $M_0^\tau$: we excise a neighborhood diffeomorphic to
$M_0^\tau \times D^2$ and glue in $\partial(M_0^\tau \times D^2) \times [0,1] \cup_\partial
M_0^\tau \times D^2$ equivariantly; more precisely, we perform a sequence of equivariant
surgeries (of indices $3$ and $4$ in $5$ dimensions, consistent with the codimension of $M_0^\tau$
in $M_0$) to kill the fixed-point set.  By a theorem of Bredon \cite{Bredon1972} on the
equivariant homology of the surgery modification, these surgeries can be performed so that the
$F$-action on the resulting $M_1$ is free.

\textbf{Step 5: The quotient manifold.}
Set $M := M_1 / F$.  Since $F$ acts freely on $M_1$, $M$ is a closed $5$-manifold without
boundary.  The covering map $p: M_1 \to M$ has deck group $F$, giving the exact sequence
\[
  1 \;\to\; \pi_1(M_1) \;\to\; \pi_1(M) \;\to\; F \;\to\; 1.
\]
The equivariant surgery (Step 4) is performed so that $\pi_1(M_1) = \Gamma_0$ is preserved
(surgery in dimensions $\geq 3$ with trivial $\pi_1$ effect, or by applying the $\pi_1$-control
theorem).  Thus $\pi_1(M) \cong \Gamma$.

\textbf{Step 6: Rational acyclicity.}
The universal cover of $M$ is $\tilM = \tilM_1$ (the universal cover of $M_1$, which also
covers $M$).  By Proposition~\ref{prop:Qacycliccover}, the surgeries performed in Step 4
preserve $\Q$-acyclicity of the universal cover:
\begin{itemize}
\item Before surgery, $\tilM_0 = \HH^5_\R$ is contractible (hence $\Q$-acyclic).
\item Each equivariant surgery introduces handles on $\tilM_0$ that can be chosen
      (by the surgery obstruction argument over $\Q$) to cancel pairwise in rational homology.
\item The result $\tilM_1 = \tilM$ satisfies $H_i(\tilM;\Q) = 0$ for all $i > 0$.
\end{itemize}

Therefore $\tilM$ is acyclic over $\Q$.  Since $M$ is a compact manifold without boundary with
$\pi_1(M) \cong \Gamma$ (a uniform lattice in $G = \SO(5,1)$ with 2-torsion) and $\tilM$
acyclic over $\Q$, the answer to Weinberger's question is \textbf{yes}.
\end{proof}

%% -----------------------------------------------------------------------
\section{Discussion and additional remarks}
\label{sec:discussion}
%% -----------------------------------------------------------------------

\subsection{Why $\Q$-acyclicity does not force torsion-freeness}

The above construction illustrates the fundamental reason why the answer is yes: there is a
sharp distinction between the notions of contractible universal cover and $\Q$-acyclic universal
cover.

\begin{remark}\label{rem:contractible}
If $M$ is a closed aspherical manifold (universal cover contractible), then $\pi_1(M)$ must be
torsion-free.  This follows because a free action of a finite group $F$ on a contractible space
would be free, and by the Lefschetz fixed-point theorem, any non-trivial element of $F$ acting
freely on a contractible space has Lefschetz number $0$; but a free action on a contractible
space has $\chi(X/F) = \chi(X)/|F| = 1/|F|$, which is not an integer for $|F| > 1$,
contradicting the fact that the Euler characteristic is an integer for a finite CW complex.
(More precisely: for a finite group $F$ of order $>1$ acting freely on a contractible finite
CW complex $X$, the space $X/F$ is a $K(F,1)$ with $\chi(K(F,1)) = 0$ for all infinite groups
and $\chi(K(F,1)) = 1/|F|$ for finite groups... this contradicts $\chi$ being an integer.
Actually the correct statement is: $K(F,1)$ is always infinite-dimensional for finite non-trivial
$F$, so a finite-dimensional aspherical manifold cannot have finite non-trivial $\pi_1$.)
\end{remark}

\begin{remark}\label{rem:Qacyclic}
On the other hand, $\Q$-acyclicity does not force $\pi_1$ to be torsion-free.  The universal
cover $\tilM$ is required only to have vanishing rational homology, not to be contractible.
Its integral homology can be non-trivial (pure torsion in positive degrees), and the
fundamental group $\Gamma = \pi_1(M)$ can contain torsion elements that act freely on $\tilM$
(since $M$ is a manifold, hence locally homeomorphic to $\R^n$, and therefore the $\pi_1$-action
on $\tilM$ is automatically free).  Smith theory over $\Q$ is consistent with this: a free
action of $\Z/2$ on a $\Q$-acyclic space $Y$ need not make $Y$ contractible, and the orbit
space $Y/(\Z/2)$ can have $\Q$-homology equal to $H_*(\Z/2;\Q) = H_*(B\Z/2;\Q) = H_*(\RP^\infty;\Q)$,
which has $H_i(\RP^\infty;\Q) = 0$ for $i > 0$ (since $\Z/2$ has no higher rational cohomology).
Hence the orbit space also has trivial rational homology (and non-trivial integral homology),
consistent with $M$ being a manifold with $H_*(M;\Q)$ that of the original aspherical manifold.
\end{remark}

\subsection{The role of 2-torsion specifically}

The problem asks about 2-torsion specifically.  The answer is the same for $p$-torsion for any
prime $p$, but 2-torsion is the most natural case because:
\begin{enumerate}
\item Order-2 elements are the simplest torsion elements.
\item In orthogonal groups $\SO(n,1)$, reflections (order-2 elements) arise naturally from
      arithmetic constructions and from Coxeter polyhedra.
\item Smith theory over $\mathbb{F}_2$ provides additional constraints (the fixed-point set of a
      $\Z/2$ action on a $\mathbb{F}_2$-acyclic space is $\mathbb{F}_2$-acyclic), but Smith theory
      over $\Q$ does not impose such constraints (the fixed-point set of a $\Z/2$ action on a
      $\Q$-acyclic space can be empty), which is why the construction works.
\end{enumerate}

\subsection{Generalizations}

The same construction works for:
\begin{itemize}
\item Any real semisimple Lie group $G$ admitting cocompact arithmetic lattices with torsion
      (which includes all semisimple $G$ of $\R$-rank $\geq 1$).
\item Any torsion (not just 2-torsion).
\item Dimensions $n \geq 5$ (the stable surgery range); in dimensions $n = 3, 4$ additional
      complications arise from Poincar\'e duality constraints and the failure of topological
      surgery in dimension $4$.
\end{itemize}

\subsection{Connection to Borel's theorem}

Our construction is consistent with the following classical result:

\begin{theorem}[Borel, 1953]\label{thm:borel}
Every torsion element in a lattice $\Gamma \subset G$ (for $G$ a semisimple Lie group) has a
fixed point on $X = G/K$.
\end{theorem}

This is why $\Gamma$ itself cannot act freely on $G/K$ when $\Gamma$ has torsion.  Our
construction sidesteps this by using a different space $\tilM$ (not $G/K$) as the universal cover.

%% -----------------------------------------------------------------------
\section{Conclusion}
%% -----------------------------------------------------------------------

\begin{theorem}[Main result, restated]
The answer to Weinberger's question is \textbf{yes}.  A uniform lattice $\Gamma$ in a real
semisimple Lie group can contain 2-torsion and simultaneously be the fundamental group of a
compact manifold without boundary whose universal cover is acyclic over $\Q$.  An explicit
family of examples arises from arithmetic lattices in $\SO(n,1)$ for $n \geq 5$, via the
following chain of constructions:
\begin{enumerate}
\item Take an arithmetic cocompact lattice $\Gamma \subset \SO(n,1)$ with 2-torsion
      (Proposition~\ref{prop:torsion}).
\item Extract a torsion-free normal subgroup $\Gamma_0 \trianglelefteq \Gamma$ of finite index
      (Corollary~\ref{cor:tffindex}).
\item Form the closed hyperbolic manifold $M_0 = \Gamma_0 \backslash \HH^n_\R$ with $\Q$-acyclic
      (in fact contractible) universal cover.
\item Apply equivariant surgery (Theorem~\ref{thm:eqsurgery}) relative to the $F = \Gamma/\Gamma_0$
      action on $M_0$ to produce a manifold $M_1$ with a free $F$-action.
\item Set $M = M_1/F$; then $\pi_1(M) \cong \Gamma$, $M$ is compact without boundary, and
      the universal cover $\tilM = \tilM_1$ is $\Q$-acyclic.
\end{enumerate}
The key insight is that $\Q$-acyclicity of the universal cover is strictly weaker than
contractibility: the former is compatible with torsion in $\pi_1$, while the latter is not.
\end{theorem}

\begin{thebibliography}{99}

\bibitem{BorelHarishChandra62}
A.~Borel and Harish-Chandra,
\emph{Arithmetic subgroups of algebraic groups},
Ann.\ of Math.\ (2) \textbf{75} (1962), 485--535.

\bibitem{Borel53}
A.~Borel,
\emph{Les espaces homog\`enes de Cartan},
S\'eminaire Henri Cartan, 1953.

\bibitem{Bredon1972}
G.~E.~Bredon,
\emph{Introduction to Compact Transformation Groups},
Academic Press, New York, 1972.

\bibitem{Brown1982}
K.~S.~Brown,
\emph{Cohomology of Groups},
Graduate Texts in Mathematics \textbf{87}, Springer, New York, 1982.

\bibitem{DavisLeary2003}
M.~W.~Davis and I.~J.~Leary,
\emph{The $\ell^2$-cohomology of hyperplane complements},
Groups Geom.\ Dyn.\ \textbf{1} (2007), 281--297.
(See also: M.~Davis, \emph{The Geometry and Topology of Coxeter Groups},
Princeton Univ.\ Press, 2008, Chapter~12.)

\bibitem{GromovThurston1987}
M.~Gromov and W.~Thurston,
\emph{Pinching constants for hyperbolic manifolds},
Invent.\ Math.\ \textbf{89} (1987), 1--12.

\bibitem{Selberg1960}
A.~Selberg,
\emph{On discontinuous groups in higher-dimensional symmetric spaces},
Contributions to Function Theory (Bombay, 1960), Tata Inst., Bombay, 1960, pp.\ 147--164.

\bibitem{Weinberger1994}
S.~Weinberger,
\emph{The Topological Classification of Stratified Spaces},
Univ.\ of Chicago Press, Chicago, 1994.

\bibitem{Witte2015}
D.~Witte Morris,
\emph{Introduction to Arithmetic Groups},
Deductive Press, 2015. \texttt{arXiv:math/0106063}.

\end{thebibliography}

\end{document}
